% =====================================================================
%  R. Nielsen, IDE OG VIRKELIGHED
%  In: For Ide og Virkelighed. Et Tidsskrift, udgivet af R. Nielsen, Bj. Bjørnson og Rud. Schmidt.  Aargang 1869.
%  Kjøbenhavn: Chr. Steen & Søns Forlag; Louis Kleins Bogtrykkeri.  pp. 1-39.  Antiqua.
%
%  TRANSCRIPTION -- GENERATED; DO NOT EDIT BY HAND.  The five Nielsen essays of the 1869 volume are built by ONE pipeline
%  (ebook-method/for-idee-1869/) into a combined working file, which split5.py cuts into the five transcriptions.
%  Rebuild from the repo root:   sh ebook-method/for-idee-1869/rebuild.sh
%
%  STATUS: COMPLETE (2026-10-13): the volume's title page and the whole essay, pp. 1-39, with its notes.
%
%  EDITION.  The only printing.  Text of record: the Internet Archive's scan of the University of Toronto copy
%  (forideogvirkelig00kjobuoft), ~/bibliotek/Nielsen, Rasmus/1869-for-idee-og-virkelighed.pdf (SCANS.tsv), 886 PDF pages,
%  400 ppi colour; text layer = the IA's ABBYY FineReader OCR.  Printed page N = PDF N + 10 to p. 688; a plate is bound in
%  before p. 689, so pp. 689-706 = PDF N + 12.  Title page PDF 7.  No e-book, no other scan, no earlier transcription found
%  (checked 2026-10-13: the repo, the web, the Internet Archive; HathiTrust and Google Books could not be queried that day).
%
%  ACCURACY (2026-10-13; one collation for all five essays, 189 pp., ~48,000 words).  The words are those of the ABBYY layer
%  (draft.py), collated word by word with an independent tesseract reading of the images (collate/check.py, witness 'ocr'):
%  1132 disagreements at the first collation.  check.py's own rules settled 96; rules.py settled ~330 more, each rule
%  switched on only after a BLIND PILOT of 107 crops (reader A, 21/22 controls) had measured it on this book: tesseract's
%  glyph faults 8/8, aa read an 8/8, accents it adds 8/8, the layer's glyph faults 7/7 and 8/8, o/ø by the ø rule 8/8, the
%  printer's space before a mark 8/8 (house rule), and the narrow word-division rules; the quote, dash, note-number and
%  word-seam rules failed the pilot and stayed off.  The rest were read at the image BLIND (sheets.py: 579 word crops on
%  58 sheets, two planted controls a sheet, 116 of 116 controls right; six readers) and about 110 spots in the main
%  session from line crops (zoom.py).  STRUCTURE read at the image (reader K): the essays' headings, the volume's title
%  page, every footnote whole (the notes were collated where they stand and then moved to their marks), the signatures.
%  LETTERSPACING (the print's emphasis, \emph): candidate lines by box width (spacedwidth.py; the gap measure is noise on
%  this scan): 163 lines read blind (22/22 roman controls), doubtful ones read in the main session.  ITALICS by stroke
%  slant (check.py --italics): 24 crops read blind; the only italics are Latin quotations.
%  SWEEPS after OPEN 0, each spot read at the image: the W-shortens / T-longer list of the blind decisions (found two
%  W decisions that had dropped a note's mark and a closing «, and one that had joined „der f.“); --moved and --splits;
%  the split sweep; the reader-line sweep; paragraph starts against tesseract's indents; quotation marks (the print
%  closes with « and often with », both as printed); the unknown-word sweep (31 spots: 16 corrections).
%  TRANSLATION as a last reading: the twelve translators listed 79 odd spots (ebook-method/for-idee-1869/tr/odd-*.tsv),
%  ~60 read at the image (zoom zo, zw, zx): mended -- f's whose head did not ink, read t by both OCRs (pp. 21, 36, 309, 319,
%  441; a recurring defective sort, as in Folkelige Foredrag), two blind decisions that had taken a reader's t for the
%  print's f (pp. 322, 539), a speck read as a stop (p. 589), a lost dash (p. 525), the verse on p. 414 set as one quotation,
%  false paragraph breaks at page turns left by the notes' removal (pp. 430, 553, 620, 695, 703); the rest are the print's.
%  SAMPLED PAGE CHECK (pagecheck.py, seed 202610136): pp. 12, 418, 521, 694 and the title page read whole and blind, 1069
%  words: 0 word and 0 punctuation errors; one letterspaced passage in a note missed (p. 694), mended.
%  What this cannot see: a misprint both OCRs normalise into a real word.
%  The ø rule (Hans 2026-10-08): o-forms that are no word (Oieblik, Ifolge, imode ...) are set with ø.
%  PRINTED AS IS p. 2: „hviken“ for hvilken
%  DEFECT, letter supplied p. 21: „fire“: the f did not ink its head (both OCRs read t); letter supplied
%  PRINTED AS IS p. 24: „Hensende“ for Henseende
%  PRINTED AS IS p. 32: „ventet. indtil“: a stop before a lower-case word
%  PRINTED AS IS p. 32: „Rette. organisere“: a stop before a lower-case word
%  DEFECT, letter supplied p. 36: „fik“: the f did not ink its head (both OCRs read t); letter supplied
%  PRINTED AS IS p. 2: „Virkelighen“ for Virkeligheden
%  PRINTED AS IS p. 4: „Polemik»“: the closing mark printed » (so also pp. 5, 22, 605, 691-694, 704)
%  PRINTED AS IS p. 22: „Sageu“: a turned n
%  PRINTED AS IS p. 32: „Resulat“
%  Not transcribed: running heads, folios, sheet signatures („For Ide og Virkelighed.  20“), the library's stamp; the large
%  initial of each essay is set as a plain letter.  Quotation marks as printed (» « in the text, „ “ in the notes).  The notes
%  are numbered ¹) per page, as printed (p. 553: *)).
%
%  TOOLS (ebook-method/for-idee-1869/, copied from folkelige-foredrag): draft.py, ocr.py, rules.py, sheets.py + reads.py +
%  READER-BRIEF.md, heads.py, spacedwidth.py + spacedlines.py + spaced.tsv, zoom.py, pagecheck.py + PAGECHECK-BRIEF.md,
%  patches.py (PATCHES, NOTES, POST), decisions.tsv, the sweeps, split5.py, rebuild.sh.
% =====================================================================
\documentclass[12pt,a4paper,openany]{book}
\usepackage[utf8]{inputenc}
\usepackage[T1]{fontenc}
\usepackage{libertinus}
\usepackage{libertinust1math}
\usepackage[danish]{babel}
\usepackage[protrusion=true,expansion=true]{microtype}
\usepackage{setspace}
\usepackage[margin=1.3in]{geometry}
\usepackage{fancyhdr}
\usepackage{hyperref}
\hypersetup{pdftitle={Ide og Virkelighed}, pdfauthor={Rasmus Nielsen}, hidelinks}
\pagestyle{fancy}
\fancyhf{}
\fancyhead[LE,RO]{\thepage}
\fancyhead[C]{\textit{R. Nielsen: Ide og Virkelighed (1869)}}
\setstretch{1.3}
\usepackage{marginnote}
\renewcommand{\thefootnote}{\arabic{footnote})}   % the print numbers each page's note ¹) (p. 553: *))
\newcommand{\opage}[1]{\marginnote{\footnotesize\textit{[#1]}}}
% the essay's heading: title (and subtitle), centred, a short rule; the print's large initial is not reproduced
\newcommand{\essayhead}[2]{\begin{center}{\large #1}\ifx&#2&\else\\[0.3em]{\normalsize #2}\fi\\[0.5em]\rule{3em}{0.4pt}\end{center}\medskip\noindent}
\newcommand{\signature}[1]{\par\smallskip\hfill #1\hspace*{2em}\par}

\begin{document}

% --- front matter: the volume's title page (PDF 7), read at the image (reader K, F7; the page check reads it again blind).
\begin{titlepage}
\begin{center}
\vspace*{3em}
{\Huge For Ide og Virkelighed.}\\[1em]
\rule{8em}{0.4pt}\\[1em]
{\large Et Tidsskrift,}\\[0.8em]
{\small udgivet}\\[0.4em]
{\footnotesize af}\\[0.6em]
{\large R. Nielsen,}\\[0.4em]
{\large Bj. Bjørnson og Rud. Schmidt.}\\[1em]
\rule{8em}{0.4pt}\\[1em]
{\large Aargang 1869.}\\[1em]
\rule{8em}{0.4pt}\\[4em]
Kjøbenhavn.\\[0.4em]
Chr. Steen \& Søns Forlag.\\[1.5em]
{\small Louis Kleins Bogtrykkeri.}
\end{center}
\end{titlepage}

% ===== TEXT (pp. 1-39) =====
\essayhead{Ide og Virkelighed.}{}
% --- p. 1 ---
\setcounter{footnote}{0}%
\opage{1}Det er en egen Sag med Ide og Virkelighed: de kunne
ikke undvære hinanden og ligge dog bestandig i Strid
med hinanden. Virkeligheden for sig uden ledende Ideer
og lysende Hovedtanker frembyder kun et bedrøveligt
Skue, et uhyggeligt Syn af delte Interesser, smaalig
Kiv, Forvirring og Opløsning; og dog er det under slige
Omstændigheder ingenlunde den vrantne Virkelighed
beleiligt at indlade sig med Ideer, allermindst med den
Ide, den dybest savner, den vækkende, klarende, foryngende
Ide, hvorefter den, uden at forstaae sig selv,
i Løndom længe har sukket.

Virkelighedens Protest mod en ny, fremtrængende
Ide er ikke noget gyldigt Beviis for Ideens Ugyldighed.
Tværtimod! En Livside kan i og for sig være fuldkommen
berettiget, kan have alle Adkomster som Prætendent
til Virkelighedens Throne og endda møde Modstand
i Omstændighederne og gjentagne Gange blive
slaaet tilbage, enten fordi det praktiske Liv endnu
staaer under Herredømmet af forældede, men ved deres
Ælde ærværdige Traditioner, der først efterhaanden
maae bortdøe, eller ogsaa fordi Virkeligheden ved pludselig
at afryste et helt Dynasti af gamle forgjemte
% --- p. 2 ---
\setcounter{footnote}{0}%
\opage{2}Ideer har faaet Blodet i Kog og er kommen i et saadant
Oprør, at der maa hengaae nogen Tid, inden den
ret kan besinde sig og skaffe Betingelserne for den
nye, livsfriske Ide tilveie. Omsider kommer den længe
forventede Stund, Forstaaelsens lykkelige Øieblik, Virkeligheden
slaaer Øiet op og kjender Ideen: Ide og
Virkelighed styrte i hinandens Favn lig to Forelskede,
der længe maatte kæmpe med Misforstaaelser. Men
lykkelige Øieblikke! desværre, de vare kun kort. Naar
Troskabseden er svoren, Indvielsens Festligheder forbi,
og Hvedebrødsdagene tilende, saa kølnes Begeistringen,
Tilværelsens Kreditorer melde sig med Regning
paa Betingelserne, saa komme Bryderierne igjen;
thi det er nu engang Virkelighedens Lod: den er altid
fuld af Bryderier. Nu følger Klage paa Klage fra begge
Sider. Ideen klager, at Virkelighen er treven, triviel   % PRINTED AS IS (p. 2): „Virkelighen“ for Virkeligheden
og aandløs, uden Alvor, uden Udholdenhed, uden Opsving.
Virkeligheden klager, at Ideen med sine vidtgaaende
Fordringer og høitflyvende Planer er upraktisk,
hensynsløs, uden Blik for det Nærmeste, det Nyttige, en
Idealist, en Fantast, rentud sagt, en Sværmer. Det kan
komme til hæftige Optrin mellem Parret; Striden kan
gaae saa vidt, at de trods gjensidige Forpligtelser og
Vielsens høitidelige Løfter endog appellere til Publikums
høistærede Retfærdighed og begjære Skilsmisse.

For at dømme i Sagen, er det hver Gang nødvendigt
at undersøge, paa hvilken Side Skylden er.
Men for at udfinde, paa hviken Side Skylden er, maa man   % PRINTED AS IS (p. 2): „hviken“ for hvilken
i ethvert paagjældende Tilfælde kunne gjøre Forskjel
mellem en sand og en falsk Ide. Er det en falsk Ide,
der har spillet Forførerens Rolle og ved blændende
% --- p. 3 ---
\setcounter{footnote}{0}%
\opage{3}Skin og smigrende Talemaader bedaaret Virkeligheden
i et ubevogtet Øieblik og narret den ind paa en farlig og
fordærvelig Vei, da er det jo ligefrem Pligt at komme
den stakkels Bedaarede til Hjælp, revse Forføreren og
see, saavidt muligt, at faae det overilede Skridt gjort
godt igjen. Men hvis Ideen er i sin gode Ret, hvis det
Maal, den har sat, er det eneste fornuftige, fordi det
er det eneste, der virkelig svarer til Tidens Tarv,
og den Vei, den har anvist, den eneste rette, fordi den
er den eneste, der fører til Maalet: saa har jo Virkeligheden
Uret, naar den, fordi det koster Anstrængelse,
vægrer sig ved at lægge Kræfterne til og ufortrøden arbeide
paa Løsningen af de Opgaver, Ideen har foreskrevet.

Hvad dette vil sige, kan sees ved et Blik paa den
Ide, som Nutiden selv har ophøiet til at være den
herskende, nemlig Frihedens Ide. Er Friheden, den
politiske Folkefrihed, naar vi undersøge Sagen upartisk
og tale uden Omsvøb, en sand, selvgyldig og berettiget
Ide eller maaskee falsk, selvmodsigende og fordærvelig?
At gjøre dette til et aabent Spørgsmaal er, i Betragtning
af de Vanskeligheder, Bryderier og Forviklinger,
hvormed den allerede har belemret Virkeligheden, vist ikke
saa ubetimeligt, som det ved første Øiekast kunde synes.
Dog nytter det ikke at begynde med Virkeligheden,
thi Virkeligheden er underkastet mange Indflydelser;
Virkeligheden har mange Sind, anlægger mange
Dragter og taler med mange Tunger. Hvor var i Aaret
1848 ikke hele den virkelige Samtid begeistret for Friheden,
beredt til at bringe ethvert Offer og vove Alt
for Ideen! Det var Opvækkelsens Tid, Forventninger%
% --- p. 4 ---
\setcounter{footnote}{0}%
\opage{4}nes gyldne Tid, Haabets Tid, og Haabet var som den
friske Vaar, Haabet var lysegrønt. Men saa ere Forventningerne
igjen siden den Tid sunkne betydelig.
Virkeligheden har allerede adskillige Gange skiftet
Farve; det Lysegrønne har undertiden tabt sig i det
Graalige og er endog paa nogle Steder gaaet over i
det Brandgule. Nei, af Virkeligheden faaer man ingen
paalidelig Forklaring; man maa begynde med Ideen.

Men er der da virkelig i vore Dage noget fornuftigt
Menneske, som for Alvor kan falde paa at opkaste
Tvivl om Folkefrihedens politiske Gyldighed? Vi svare
Ja! Der er endnu adskillige ret fornuftige Mennesker,
der nære nogen Tvivl. Vi skulle til Bevis for vor Paastands
Rigtighed anføre en i denne Henseende original
Tvivler, en Tænker af første Rang, en klartseende, frisindet
Mand, en Mand, som har Samtidens enstemmige
Vidnesbyrd for, at han idetmindste har holdt sig udenfor
alle politiske Partier, og forsaavidt maa kaldes uinteresseret,
en Mand af overlegen Kløgt og med sikkert Blik for
Virkelighedens Bryderier, en Mand, der netop satte det
som sit Livs Opgave at gjennemtænke Frihedens høieste
Ide, og saa, mærkværdigt nok, kom --- som det synes
--- til det Resultat, at Folkefrihed, politisk forstaaet,
er en selvisk Indbildning, en falsk Ide. Vi mene S.
Kierkegaard.

S. Kierkegaard var, som bekjendt, ingen Ynder af
Tidens Fordringer, ingen Beundrer af de saakaldte liberale
Bestræbelser. Allerede hans første Artikler om »Kjøbenhavnspostens
Morgenbetragtninger« og »Fædrelandets
Polemik» samt de Repliker han i den Anledning vexlede   % PRINTED AS IS (p. 4): „Polemik»“: the closing mark printed » (so also pp. 5, 22, 605, 691-694, 704)
med »Hr. Orla Lehmann«, ere betegnende. Det var en
% --- p. 5 ---
\setcounter{footnote}{0}%
\opage{5}Stemningspolemik, hvori han til en Begyndelse prøvede
sin spidse Pen; det var hans første Greb paa Idealitetens
Strængeleg, et Par raske Buestrøg, for at høre Tonen,
og Tonen angav tydelig nok, at denne Violin ikke vilde
stemme med Politikens blæsende Instrumenter.

Paa et senere Stadium, hvor Talen særlig er om
Stat og Kirke, udtaler han sig, foranlediget ved en
Yttring af Dr. Rudelbach, meget skarpt om hvad han
kalder en usalig Forvexling af Politik og Kristendom,
den bagvendte Reformation, der reformerende sætter et
nyt Slettere i Stedet for et gammelt Bedre, »medens det
dog skal være vist og sandt, at det er en Reformation,
i hvilken Anledning hele Byen bliver illumineret.« Den
sande Frihed er ikke en udvortes Foranstaltning, men
en indvortes Sag; Frihed er væsenlig Inderlighed. »Som
det er Menneskets Fortrin for Dyret at kunne leve
under ethvert Klima, saaledes rangerer Fuldkommenhed
i Kristendom, just fordi denne er Inderlighed, efter sin
Styrke til at kunne leve, om saa var, under de ufuldkomneste
Former. Politik er dette udvortes Væsen,
denne tantaliske Travlhed for Forandring i det Udvortes».

Frihedserkjendelsen er i strengeste Forstand Selverkjendelse;
Frigjørelsen kan derfor ikke begynde med
noget Udvortes, ikke tage sit Udgangspunkt hverken i
Slægten eller i Verdenshistorien, hverken i nye Stats- eller
i nye Kirkeforfatninger; den maa begynde indenfra,
idet hver Enkelt begynder med sig selv, ikke med
store Forventninger om »de frie Institutioner« , men i
Tro paa Sandhedens frigjørende Aand.

Det er et Sandsebedrag, naar Nogen mener, at det
er de frie Institutioner, der skal frelse Kirken og Kri%
% --- p. 6 ---
\setcounter{footnote}{0}%
\opage{6}stendommen. »Skulde dette«, siger K., »høre med til
sand Kristendom, denne Tro paa de, politisk opnaaede,
frie Institutioners frelsende Magt, saa er jeg ingen
Kristen, ja endnu værre, jeg er et rent Djævelens Barn;
thi oprigtig talt, jeg har endog Mistanke til disse, politisk
opnaaede, frie Institutioner, især til deres frelsende
og gjenfødende Magt. Det er min Kristendom,
eller saa kristendum er jeg, som da forresten aldrig
har befattet mig med »Kirke« og »Stat« --- det er meget
for stort for mig; dertil behøves ganske anderledes
Seere end jeg er, eller ogsaa bør det ganske overlades
de ordenligen dertil Beskikkede og Kaldede. Jeg
har aldrig kæmpet for »Kirkens« Emancipation saa lidet
som for den grønlandske Handels, Kvindernes, Jødernes
eller nogensomhelst anden Emancipation. Jeg har kun som
Enkelt, med Sigtet »den Enkelte«, tilsigtende Inderliggjørelse
i Kristendom i »den Enkelte«, konsekvent med
Aandens Vaaben, ene og alene med Aandens Vaaben
kæmpet for at gjøre opmærksom paa og mod at lade
sig bedrage af Sandsebedrag«.

Som det da er et Sandsebedrag, om Nogen mener,
at det er de ydre Former, der hindre ham i at naae
Friheden, saa er det jo ogsaa et Sandsebedrag, ja det
Samme vendt paa en anden Maade, om Nogen bilder
sig ind, at det er de ydre Former, der skulle forhjælpe
ham til Friheden.

Ifølge denne Synsmaade kunde S. K. da heller ikke
knytte store Forhaabninger til de frie Bevægelser, der
udmærkede Aaret 1848. Hans lille Skrift: »Om min
Forfatterviksomhed« (1851) ender med følgende, for hele
hans Standpunkt mærkelige Stykke: »De brast Aa%
% --- p. 7 ---
\setcounter{footnote}{0}%
\opage{7}ret 48, Klogskabens Traade; det hørtes det Hvin,
der bebuder Chaos! »Det var det Aar 48, det var
et Fremskridt«, ja --- hvis der nemlig skaffes Regering
tilveie, hvortil maaske ikke behøves en eneste
ny Embedsmand, eller nogen ældres Afskedigelse, men
vel maaske en inderlig Forvandling i Retning af at blive
fast ved Gud. Feilen fra oven var vistnok den, at i det
Hele taget fra øverst til nederst Regeringen igjennem
var Styrken væsenligen verdslig Klogskab, hvilket just
væsenligen er Mangel paa Styrke; Skylden fra neden
var den, at ville have al Regering bort: Straffen ---
thi Syndens Maade er altid Straffens Maade --- Straffen
er: det, der bittrest bliver savnet, er just Regering.
Aldrig som i vort Aarhundrede var vel Slægten og de
Enkelte i den (den Bydende --- den Lydende; den Foresatte
--- den Underordnede; Læreren --- den Lærende
o. s. v.) saaledes frigjorte fra, om man saa vil, alle
Gener af, at Noget staaer og skal staae ubetinget fast;
aldrig følte vel »Meningerne« (de mest uensartede,
paa de forskjelligste Gebeter) sig saaledes i »Frihed, Lighed
og Broderskab«, ugenerte og lyksaliggjorte under
Frispasens »til en vis Grad«: aldrig vil vel Slægten saa
dybt komme til at sande, at den og hver Enkelt i den
trænger til, at Noget skal staae ubetinget fast, trænger
til, hvad Guddommen, den Kjærlige og i Kjærlighed,
opfandt, det Ubetingede, istedenfor hvilket Mennesket,
den Kloge til egen Fordærvelse, beundrende sig selv satte
dette Beundrede »til en vis Grad«. Byd Sømanden at
seile uden Baglast --- han kulseiler; lad Slægten, lad
hver Enkelt i den forsøge at bestaae uden det Ubetingede:
det er og bliver Hvirvel. Det kan i længere eller
% --- p. 8 ---
\setcounter{footnote}{0}%
\opage{8}kortere Tid, i Mellemrum synes anderledes, at det er
Fasthed og Tryghed, i Grunden er og bliver det Hvirvel,
selv de største Begivenheder og det mest anstrængede
Liv dog Hvirvel eller som det at sye uden at fæste
Ende --- indtil Enden atter gjøres fast derved, at det
Ubetingede anbringes, eller derved, at den Enkelte, om
end paa nok saa stor Afstand, dog forholder sig til et
Ubetinget. Leve blot i det Ubetingede, indaande kun
det Ubetingede, kan Mennesket ikke, han omkommer
som Fisken, naar den skal leve i Luften; men paa den
anden Side, uden at forholde sig til det Ubetingede kan
et Menneske heller ikke i dybere Forstand »leve«, han
udaander, det er, han vedbliver maaske at leve, men
aandløst. Er --- for at blive ved Mit: det Religiøse
--- er Slægten eller en stor Del Enkelte i Slægten
voxede fra dette Barnlige, at et andet Menneske for
dem er Den, der repræsenterer det Ubetingede: vel,
derfor kan alligevel det Ubetingede ikke undværes. Saa
maae de Enkelte selv forholde sig til det Ubetingede.
Det er dette jeg efter forundt Evne, med yderste Anstrængelse
og mangen Opoffrelse har kæmpet for, kæmpende
mod ethvert Tyranni, ogsaa det Numeriskes.
Denne min Stræben er bleven lagt for Had som uhyre
Stolthed og Hovmod --- jeg troede, og jeg troer, at dette
er Kristendom og Kjærlighed til Næsten.«

Hvorledes en fri Forfatning end ordnes --- her tales
om frie Forfatninger i al Almindelighed --- maa det tilsidst
gaae ud baade over Stat og Borgere. Regeringsmyndigheden
svækkes. Ingen Enkelt enten tør eller kan
træde afgjørende op, handle af egen Overbevisning og
tage Ansvaret paa sig; men Alt betinges af Øieblikkets
% --- p. 9 ---
\setcounter{footnote}{0}%
\opage{9}Stemning, Alt gjøres afhængigt af Flertallets Samtykke.
Skaden er dobbelt. For det Første maa den Sagkyndige
finde sig i, at hans egenlige Hovedtanke, hans omhyggelig
gjennemarbeidede Plan under Oppositioner og Diskussioner
beklippes, beskæres, og ved haardnakkede Indvendinger,
delvise Ændringer imødekommende Tillæmpninger
saaledes omkalfatres, at den langtfra at blive af een Støbning,
tvertimod bliver et Flikværk af uensartede Stykker,
en broget Sammensætning af ulige Mønstre, mindre
lig en Plan end et Lappeteppe. Naar et saadant Flikværk
er sat igjennem, og den bittre Erfaring langt om
længe har viist, hvor ubrugeligt det er: hvis er saa
Skylden? Ingens! Hvem bærer Ansvaret? Ingen! Dette
er endnu det allerværste. Ansvarligheden fordamper i
Diskussioner, bortdunster i Komitébetænkninger og Afstemninger,
saa at det Hele tilsidst kommer ud som
Værk af en Maskine, et Slags anonymt Produkt fra
Partiernes Drivhjul. Forsøger man at sætte et Bedre
i Stedet for det anerkjendte Slettere, saa begynder den
politiske Komedie forfra. Alt foregaaer offenlig og er
dog et fortsat Intriguespil; det seer ud, som om hver
enkelt Sag blev underkastet en alsidig Prøvelse, en omhyggelig, vidtløftig Begrundelse; men hvad er det, der
egenlig prøves? Det er ikke Tingen selv, men Partiernes
Kræfter. Hvad giver Grundene Vægt? Ikke Sagens
Natur, ikke overbevisende Tanker, men de Stemmegivendes
helt eller halvt maskerede Tendenser. Det i og
for sig Sande, Fornuftige, Gavnlige, kan naturligvis ikke
uden videre komme i Betragtning; thi Betingelserne ere
langt vigtigere end det Ubetingede. Det gaaer slet ikke
an at rykke ud med Sproget, forsvare det, efter Ens
% --- p. 10 ---
\setcounter{footnote}{0}%
\opage{10}egen Overbevisning, eneste Rette, overladende til Omgivelsen,
hvad den vil antage eller forkaste, nei det gjælder
om at fange Vinden, strække Følehornene ud, sondere
Partierne, for at erfare hvad man kan sætte igjennem. At faae Noget sat igjennem er det Vigtigste; det
Næstvigtige er, hvad det er for Noget: en halv Sandhed,
naar den blot kan gaae igjennem, er bedre end en hel,
der udsætter sig for at blive afvist. Paa denne Maade
bliver Alt antageligt og Alt forkasteligt til en vis Grad,
Sandhed til en vis Grad, og som Følge heraf ogsaa
Usandhed, Retfærdighed til en vis Grad, følgelig ogsaa
Uretfærdighed, Redelighed til en vis Grad, følgelig ogsaa
Uredelighed. At være ligefrem usand, uretfærdig, uredelig
duer ikke; men at være ubetinget sand, ubetinget
retfærdig, aaben og redelig duer heller ikke: Virkeligheden
kan ligesaa lidt bruge den rene Sandhed som den
rene Usandhed; den kræver en egen Blanding af Godt
og Slet, af Redelighed og Uredelighed. Virkelighedens
Karakter bliver i denne Forstand en klog Forfarenhed,
en kløgtig beregnet Karakterløshed. Men den kløgtige Karakterløshed
er sædelig Tvetydighed, en raffineret Underfundighed,
der ikke aabenlyst angriber, men listig
udhuler det Bestaaende. Dette er for den borgerlige Tænkemaade
i det Hele det Dyssende, Bedøvende, Demoraliserende.
»Thi mod Oprør kan man bruge Magt,
for vitterligt Falskmaal er Straffen i Vente; men dialektisk
Hemmelighedsfuldhed er vanskelig at udrydde;
der skal allerede forholdsvis et finere Øre til for at spore
Reflexionens lydløse Listen ad Tvetydighedens Snigveie.
Det Bestaaende bestaaer, men Reflexionens Lidenskabsløshed
finder sin Beroligelse i, at det er en Tvetydighed.
% --- p. 11 ---
\setcounter{footnote}{0}%
\opage{11}Man vil ikke have Kongemagten afskaffet, ingenlunde,
men hvis man lidt efter lidt kunde faae den forvandlet
til en Indbildning, saa vilde man med Fornøielse raabe
Hurra for Kongen. Man vil ikke have Udmærkelsen
styrtet, men hvis man paa samme Tid kunde avancere
en Viden om, at den er en Indbildning --- saa vilde
man beundre. Man vil lade hele den kristelige Terminologi
bestaae, men underhaanden være vidende om,
at der ikke skal tænkes noget Afgjørende ved den. Og
man vil være angerløs; thi man nedriver jo Intet«.

Til de frie Forfatningers mørkeste Skyggepartier hører
just det, at Folkevillien altid middelbart eller umiddelbart
maa give sig tilkjende som Mængdens Villie, og
Mængde er, det er Ks. Thesis, »Mængde er Usandheden.«
Der er i Grunden Ingen, der har mere Foragt for Mennesket
som Menneske end de, »der gjøre Profession af at staae
i Spidsen for Mængde. Lad En komme til en Saadan,
et enkelt Menneske --- jo vist, hvad bryder han sig
derom; det er meget for lidt; stolt afviser han ham;
der maa i det mindste være Hundrede. Og naar der er
Tusinde, da bøier han sig for Mængden, bukker og skraber;
hvilken Usandhed! --- At Mængden som Mængde er
Alt, den Enkelte i Mængden Intet, betyder, at Mængden
er Publikum. Men just den Omstændighed, at den
Enkelte i Mængden er Intet, gjør det muligt, at han
kan identificere sig med Mængden, tale og skrive i Mængdens
Navn uden at staae personlig til Regnskab for
hvad han siger. Thi hvem er det, man har for sig?
Mængden og dog Ingen, en Anonym. »At en Anonym,
ved Hjelp af Pressen, Dag ud og Dag ind kan faae
(ogsaa i Henseende til det Intellektuelle, det Ethiske,
% --- p. 12 ---
\setcounter{footnote}{0}%
\opage{12}det Religiøse) sagt, hvad han vil, hvad han maaske
ikke i fjerneste Maade havde Mod til at sige end det
Mindste af i Enkelthedens Situation; kan hver Gang
han lukker sin --- ja Mund kan man da ikke kalde det
--- sit Svælg op, \emph{paa een Gang} henvende sig til Tusinde
Gange Tusinde; kan faae Tusinde Gange Tusinde
til at sige det Sagte efter --- og Ingen har Ansvaret,«
dette anseer K. for vor Tids store Elendighed, en Elendighed,
der er større end Oldtidens største. »Man sige
ikke, at saa kan jo Sandheden igjen ved Hjælp af Pressen
indhente Løgnen og Vildfarelsen. Du, der taler
saaledes, spørg Dig selv: tør Du paastaae, at Menneskene
i Mængde taget, ere ligesaa hurtige til at gribe
efter Sandhed, der ikke er velsmagende, som efter
Usandhed, der altid er lækkert tillavet, end sige, naar
det ovenikjøbet skal være forbundet med en Vedgaaen
af, at man har ladet sig bedrage! Eller tør Du blot
paastaae, at Sandheden lige saa hurtig lader sig forstaae
som Usandhed, der ingen Forkundskaber fordrer,
ingen Skole, ingen Tugt, ingen Afholdenhed, ingen Selvfornægtelse, ingen redelig Selvbekymring, ingen langsom
Arbeiden! Nei, Sandheden, der ogsaa afskyer denne
Usandhed: som eneste Formaal at eftertragte Udbredelse,
er ikke saa hurtig tilbens.« (»Synspunktet for min
Forfattervirksomhed.« S. 95—96.)

Det mørke Billede fuldender sig i Skildringen af Publikum.
Vistnok kan man ingenlunde sige, at det er Folkefrihedens
Ide, der frembringer Publikum; thi det er Kulturen:
Nutidens Publikum er en eiendommelig Udartning
af Nutidens Kultur. Derimod kan man med Grund sige,
at den frie Forfatning har sin tvetydigste Ven, sin far%
% --- p. 13 ---
\setcounter{footnote}{0}%
\opage{13}ligste Forbundsfælle i Publikum. Hvad er da Publikum?
J. L. Heiberg, der, som bekjendt, har leveret adskillige
varme Batailler med Publikum, har tillige leveret værdifulde
Bidrag til dets Signalement. I sin Afhandling »Folk og
Publikum« (Intelligensblade, 1842, Nr. 6) gaaer han ud
fra den bekjendte Modsætning, at et Folk er en Organisme,
et Publikum en Masse, og udtaler sig om Overgangen
fra Folk til Publikum saaledes: »Folkets sukcessive
Forvandling til Publikum er et Faktum, som hertillands
allerede har sin \emph{Historie}, og jeg selv har leveret
ikke faa Aktstykker dertil, skjøndt tildels uden at
være mig det bevidst, thi da vor politiske Udvikling er
meget yngre end vor literære, men det netop er hiin,
som først kaster det rette Lys paa denne, saa kunde
man længe være uvidende om den sande Betydning af
de i Literaturen sig bevægende Kræfter. Allerede i
»Julespøg og Nytaarsløier« indførte jeg Publikum personificeret
som »en kolossal Vanskabning«, men lod dog
Digteren i første Øieblik hilse den med de Ord: »Mit
Folk!« Først siden, da han ved Synet af dets Tusinder
af Øine, Munde og Hænder, og ved dets raae, sig selv
modsigende Yttringer faaer Syn paa dets atomistiske
Natur, begegner han det med den yderligste Ringeagt,
uden dog at tage det Udtryk tilbage, hvormed han først
hilste det«.

Det skjønnes nu let, at den offenlige Mening, der
spiller en saa stor Rolle i alle frie Statssamfund, kan
under Skin af, at det er Folkets Stemme, tages i Beslag
af Publikum. Just fordi Publikum altid er Tilskuer og
den Enkelte i Publikum væsenlig Ingenting --- forsaavidt
han ellers er Noget, maaske endog noget Dygtigt,
% --- p. 14 ---
\setcounter{footnote}{0}%
\opage{14}træder han strax udenfor Publikum --- saa er Publikum
ikke handlende; dets eneste Handling er at give sin
Mening tilkjende i Bifalds- og Mishagsyttringer. Tilkjendegivelsen
selv er uartikuleret. Det er en Snakken
og Praten for og imod, en Summen og Brummen, Surren
og Snurren, hvortil en Politiker maa spidse Øren,
hver Gang han vil udlytte, om det er Bifald eller Mishag,
der egenlig har Overvægt. Hvad der fremkogler
Fantomet af en offenlig Mening, er ingen hel, enfoldig
Tanke, men en yderlig sammensat Størrelse, en Sværm
af talløse smaa, infusoriske Meninger, en Sum af Brøker
med forskjellige Nævnere og forskjelligt Fortegn.
At en saadan kollektiv Størrelse umulig kan give en
Anskuelse, hvori der er folkelig Enhed, endnu mindre
være bestemmende for en sand Folkevillie, er indlysende.
»Kun i en lidenskabsløs, men reflekteret Tid« --- siger
S. Kierkegaard, som i denne Henseende ærlig fortsatte,
hvad J. L. Heiberg havde begyndt --- »kan det Fantom,
man kalder Publikum, udvikle sig ved Hjælp af Pressen,
naar denne selv bliver en Abstraktion. I begeistrede
Tider, i lidenskabelige tumultuariske Tider, selv naar et
Folk vil realisere den ufrugtbare Ørkens Ide, at sløife
Alt: det er dog intet Publikum\dots{} Publikum er et
Begreb, som slet ikke kan forekomme i Oldtiden, fordi
Folket selv en masse in corpore maatte træde op i
Handlingens Situation, maatte bære Ansvaret for, hvad
den Enkelte af dets Midte afstedkom, medens igjen den
Enkelte --- personligen som denne Bestemte maatte være
tilstede, maatte underkaste sig Øieblikkets Standret i
Bifald eller Misbilligelse.«

At Publikum ikke kan være Folket selv, ja end
% --- p. 15 ---
\setcounter{footnote}{0}%
\opage{15}ikke forestille Folket, er tydeligt blandt Andet deraf,
at Publikum ikke kan handle. I det Øieblik Publikum
vil optræde handlende, ophører det at være Publikum;
man kan da efter Omstændighederne kalde det Masse,
Mængde, Pøbel, Folk, hvad man vil, men Publikum er
det ikke. Hvor der er Tale om en Revolution hedder
det da heller aldrig, at der er skudt paa Publikum, men
at der er skudt paa Massen, paa Folket: paa et Fantom
kan man ikke skyde. Heller ikke vil Historien nogensinde
kunne melde om, at et Publikum har enten
styrtet en Throne eller jaget en Fyrste i Landflygtighed,
man kunde da hellere antage, at det var Spøgeri, der
havde styrtet Thronen, at et Spøgelse havde forjaget
Fyrsten. I Forhold til Alt, hvad der hører under Handling
og Karakter, er Publikum et Spøgelse. Publikums
Karakter er den reflekterende, ræsonnerende Ørkesløshed,
dets vexlende Meninger ørkesløse. Ja saa spøgelseagtig
ørkesløse ere disse Meninger, at det endog, hvad S. K.
træffende har viist, er ørkesløst at ville være af samme
Mening som Publikum. »At antage samme Mening som
disse og disse bestemte Mennesker, betyder, at man
veed, at disse ville være underkastede samme Farer som
En selv, at de ville fare vild med En, hvis Meningen
er vildfarende o. s. v. Men at antage samme Mening
som Publikum er en svigefuld Trøst, thi Publikum er
kun til in abstracto. Medens derfor ingen Majoritet
nogensinde har været saa sikker paa at beholde Retten
og Seiren som Publikum er det, saa er dette kun lidet
trøsteligt for den Enkelte, thi Publikum er et Fantom,
som ikke tillader nogen personlig Tilnærmelse. Dersom
% --- p. 16 ---
\setcounter{footnote}{0}%
\opage{16}En idag antager Publikums Mening og imorgen bliver
udpeben, saa bliver han udpeben af Publikum«.

Publikum kan Intet organisere, det kan kun nivellere:
Nivelleringen er »den sig etablerende Misundelse«.
Nivelleringen er ikke en Enkelts Handling, men et Reflexionens
Spil i en abstrakt Magts Haand; denne abstrakte
Magt er Publikum. »Pressens Abstraktion (thi
et Blad, en Avis er ingen statsborgerlig Konkretion og
kun i abstrakt Forstand et Individuum) afføder Abstraktionens
Fantom: Publikum, der er det egenlig Nivellerende\dots{}.
Jo mindre Ide der er i en Tid, jo mere
den, afmattet ved en opblussende Begeistring, udhviler sig
i Indolens, hvis vi endog vilde tænke os, at Pressen blev
svagere og svagere, fordi ingen Begivenhed, ingen Ide
griber Tiden: desto lettere vil Nivelleringen blive en
fordærvelig Lyst, en Sandsepirring, som et Øieblik kildrer
og kun gjør det Onde værre. Og har man ofte
skildret Eneherredømmets Demoralisation, revolutionære
Tiders Forfald, saa er en lidenskabsløs Tidsalders Forfald
noget ligesaa Fordærveligt, om end formedelst Tvetydigheden
mindre paafaldende. Denne dvaske Mængde,
som Intet selv forstaaer, og Intet selv vil gjøre, dette
Galleri-Publikum søger nu Tidsfordriv og hengiver sig
da til den Indbildning, at Alt, hvad Nogen gjør, skeer
for at det kan faae Noget at snakke om. Dvaskheden
sidder fornemt med Benene overkors, og Enhver, der vil
arbeide, Kongen og Embedsmanden og Folkets Lærer
og den dygtigere Journalist og Digteren og Kunstneren,
Alle blive ligesom spændte for, for at slæbe denne
Dvaskhed frem, der fornemt troer, at de Andre ere Hestene.
--- Saa holder da Publikum sig en Hund for Mor%
% --- p. 17 ---
\setcounter{footnote}{0}%
\opage{17}skabs Skyld. Denne Hund er den literære Foragtelighed.
Viser der sig nu en Bedre, maaske endog en Udmærket,
saa hidses Hunden paa, og Morskaben begynder.
Den bidske Hund river ham i Kjoleskøderne, tillader sig
alle uvorne Kaadheder --- indtil Publikum bliver kjed
deraf og siger: nu kan det være nok. Da har Publikum
nivelleret. Den Bedre, den Stærkere er mishandlet, og
Hunden, ja den bliver ved at være en Hund, som Publikum
selv foragter\dots{}. Og Publikum er angerløst, thi det
er jo ikke Publikum --- det var jo Hunden; ligesom
man siger til Børn: det var Katten. Og Publikum er
angerløst, thi det var jo ingen Forkleinelse --- det var
blot lidt Morskab. Og Publikum vil være angerløst, thi
det er jo egenlig ikke dem, som holde Hunden, de subskribere
blot, de hidse den heller ikke direkte paa, de
fløite heller ikke ligefrem ad den; i Tilfælde af Søgsmaal
vilde Publikum sige: Hunden er ikke min, det er
en herreløs Hund; og i Tilfælde af, at Hunden blev
sendt ud paa Veterinærskolen for at slaaes ihjel, kunde
Publikum endnu sige: det var virkeligen godt, at den
slemme Hund blev dræbt, det har vi alle ønsket --- endogsaa
Subskribenterne.« (En literær Anmeldelse. S. 95-96.)

Resultatet er let at uddrage: Folkefriheden betinges
af, at enhver ubetinget Autoritet til en vis Grad svækkes
og nedbrydes, saa at al Myndighed, ogsaa Regeringsmyndigheden,
lidt efter lidt synker ned under Relativitetens
Niveau. Gjennem seige, men bøielige, i alle
Retninger bevægelige Forfatningsformer arbeide Partierne
med deres Førere under bedst mulig Benyttelse af Svingningerne
i den offenlige Mening paa at fremme de Interesser,
som, hvad enten de for Resten ere Samfundet
% --- p. 18 ---
\setcounter{footnote}{0}%
\opage{18}skadelige eller gavnlige, fornuftige eller ufornuftige, til
en given Tid vise sig herskende hos Flertallet. Flertallets
Mening stemples jo rigtignok officielt til at gjælde
for Folkets Mening, men burde dog, eftersom den hver
Gang maa søge sit egenlige Tyngdepunkt i det overveiende
Flertal, altsaa i Mængden, allerhelst opfattes
som et forfatningsmæssigt, skjøndt midlertidigt Udtryk
for Mængdens Mening. Saa følge Metamorfoserne som
af sig selv: Regeringen bliver et Folk, Folket bliver en
Mængde, Mængden bliver et Publikum, Publikum bliver
et Fantom, et politisk Spøgelse, en Slags ubestemmeligt
Ingenting. Det er tilbageskridende Metamorfoser, Præmisser,
hvoraf der sluttes: ergo er Folkefriheden en
forfængelig Indbildning, en falsk Ide.

Dersom nu denne Slutning var fuldkommen paalidelig
og det vundne Resultat i Et og Alt uomstødeligt:
da vilde det tilvisse see sørgeligt ud i Nutiden, ikke
mindst hos os. Men Slutningen er, hvad den opmærksomme
Læser vist allerede har sagt sig selv, ikke paalidelig.
Der er i Præmissernes Tvetydighed en Sammensætning
af uligeartede Elementer, en betænkelig
Blanding, som først kan skilles og klares, naar vi besinde
os paa Forholdet mellem Ide og Virkelighed.

De, som under uhyggelige Indtryk af, hvad der
støder dem i Virkeligheden, have vænnet sig til at indrette
deres Domme efter deres Stemning og desaarsag
mene, at de under Folkefrihedens Tryk virkelig have
Grund til at være forstemte, ville hos en Forfatter som
S. Kierkegaard finde Udtalelser, der svare til deres personlige
Anskuelse, lindrende Indfald, polemisk husvalende
Udfald, Yttringer saa træffende, at det stundom maa
% --- p. 19 ---
\setcounter{footnote}{0}%
\opage{19}forekomme dem, som havde han læst i deres politiske
Hjerte og taget hvert Ord ud af deres eget Indre; de
kunne forsaavidt læse ham med aandsaristokratisk Selvfølelse
til ironisk Opbyggelse og more sig selv med de
vittige Billeder; --- men erklære ham for Folkefrihedens afgjorte
Modstander og ligefrem tage ham til Indtægt kunne
de ikke. Med ham er det en egen Sag; hos ham betyder
Angreb væsenlig Forsvar, og Forsvar igjen Angreb.
At ville holde ham fast i en enkelt ensidig Stemning,
forvexlende hine glimrende, vulkanagtige Udbrud
af idesprudende Stemningspolemik med konservative
Jerimiader af politisk Forstemthed, er uforstandigt; thi
vel er han stemningsrig, men han er aldrig forstemt;
der er i hans Grundstemning et sydende Væld af tusindfarvede
Idestemninger: han er en Mand med mere end
eet Ansigt, en bifrontisk Skikkelse.

Der er Ingen hos os, der har udtalt sig saa strengt
konservativt og dog saa yderlig revolutionært, Ingen, der
i den Grad har forsvaret og i den Grad angrebet det
Bestaaende, som S. Kierkegaard. Hvilken Modsætning,
hvilken himmelraabende Forskjel imellem hans første og
hans sidste Udtalelse om Forholdet, vel at mærke Grundforholdet,
mellem Religion og Stat! At han ikke var
politisk, men religiøs Forfatter, gjør i denne Henseende
Intet til Sagen.

I ovenanførte Artikel mod Dr. Rudelbach har han
saa godt som Intet at indvende mod den bestaaende
Kirkeforfatning, Alt er indtil videre godt, saare godt,
som det er. Det er ikke Statens Protektion, ikke
de kirkelige Indretninger, men kun en Smule verdslig
Letsindighed, der er Skyld i, at Kristendommen engang
% --- p. 20 ---
\setcounter{footnote}{0}%
\opage{20}imellem synker ned til, hvad Dr. R. kalder \emph{Vane-Kristendom.}
»Jeg hader,« siger S. K., »Vane-Kristendommen
i hvilkensomhelst Skikkelse den viser sig.«
Dette »i hvilkensomhelst Skikkelse« kunde jeg ønske
særligen bemærket; thi Vane-Kristendommen kan jo
have mange Skikkelser. Og var der slet intet Andet at
vælge, var Valget kun mellem denne Art Vane-Kristendom:
en verdslig Letsindighed, der sorgløst lever hen i
den Indbildning at være Kristen, uden maaske engang
at have noget Indtryk af Kristendommen, og den Art
Vanekristendom, der findes hos Sekterere, Opvakte,
Hyperorthodoxe, Partigængere naar galt skal være,
jeg vælger ubetinget den første. Den første Art har dog
kun letsindigt og negativt taget Kristendommen forfængeligt,
hvis man ellers overhovedet altid endog tør dømme
den saa strengt; den anden har maaske i aandeligt
Hovmod, i hvert Tilfælde positivt taget Kristendommen
forfængeligt; den første kunde man næsten fristes til at
smile af, fordi man haaber; den anden kan bringe En
til at gyse.« Her er ikke Tanke om, at der kunde være
noget Misligt i hele det statskirkelige Væsen, ikke fjerneste
Antydning af, at Præsternes Lønning som kongelige
Embedsmænd skulde staae i Strid med deres Kald
som Ordets Tjenere. Men hvorlænge varede det vel,
inden Bladet vendte sig! Fire Aar efter læste vi i
»Øieblikket Nr. 3, 6« følgende Bekjendtgjørelse. »Saalænge
der i Danmark existerer 1000 kongelige Levebrød
for Lærere i Kristendom, er der gjort det bedst Mulige
for at forhindre Kristendom. Saalænge der er 1000 kongelige
Levebrød, vil der bestandig være et dertil svarende
Tal af Mennesker, som agte paa den Maade at
% --- p. 21 ---
\setcounter{footnote}{0}%
\opage{21}fortjene deres Brød. Blandt disse vil der være nogle
Faa, som dog maaske fandt et Kald til at forkynde Kristendom.
Men just i det Øieblik, da det ret skulde
blive Alvor, ved ene i Tillid til Gud for egen Risiko at
maatte overtage at optræde som Lærer --- just da aabner
Staten dem den Magelighed, at tage et kongeligt
Embede, hvorved disse Faa, kristeligt, forkluddres. Det
langt større Antal vil slet ikke have noget Kald til at
forkynde Kristendom, men betragte det simplement som
Levebrød. Paa den Maade naaer Staten at fylde hele
Landet med bedærvet Kristendom, den allerstørste Vanskelighed
for Anbringelsen af sand Kristendom, langt
større end fuldkommen Hedenskab«.

Ingen, der ret vil fæste sin Opmærksomhed paa
disse høist uligeartede Udtalelser, kan, hvor megen Agtelse
og Sympathi, han forresten nærer for den overordenlig
begavede Forfatter, dølge for sig selv, at her
er en Modsigelse. Her handles jo ikke om forskjellige
bedre eller slettere Tilstande, thi en saadan Vexlen angaaer
kun Virkeligheden; nei, her handles om Ideen som
Ide, om et for alle Tilstande gjældende Princip. Kunde
det i Aaret 1851 stemme med Kristendommens Væsen,
at dens ordenlig beskikkede Lærere lønnedes af Staten,
saa maatte det Samme jo ogsaa gjælde for Aaret
1855, og omvendt. Enten --- eller: enten maa Kristendommen
i Løbet af fire Aar have skiftet Princip, eller   % PRINTER'S DEFECT (p. 21): „fire“: the f did not ink its head (both OCRs read t); letter supplied
ogsaa maa Forfatteren i det Mellemrum ganske have
forandret sin Dom om Kristendommens Væsen; da nu
det Første er utænkeligt, og det Sidste usandsynligt,
saa paatrænger sig uvilkaarlig det Spørgsmaal: hvorledes
skal Modsigelsen løses?

% --- p. 22 ---
\setcounter{footnote}{0}%
\opage{22}Til det nye Oplag af »Indøvelse i Kristendom«, der
udkom midt under Striden, føiede Forfatteren (»Fædrelandet«
16de Mai 1855) jo rigtignok nogle i denne Henseende
opklarende Bemærkninger. »Denne Bog, har jeg,»
siger han, »ladet udgaae i et aldeles \emph{uforandret} Oplag,
fordi jeg betragter den som et historisk Aktstykke.
Hvis den nu skulde være udgaaet, nu, da jeg har forvisset
mig om, at det Bestaaende, kristeligt, er uholdbart,
vilde den være bleven forandret\dots{}. Min tidligere
Tanke var: skal det Bestaaende kunne forsvares, er
dette den eneste Maade: ved digterisk (derfor af en
Pseudonym) at anbringe Dommen over det, ved saa at
trække i anden Potens paa »Naaden«, saa det blev Kristendom,
ikke blot ved »Naaden« at finde Tilgivelse for
det Forbigangne, men ved Naaden at finde et Slags Aflad
for den egenlige Kristi Efterfølgelse og den egenlige
Anstrængelse ved at være Kristen. Paa den Maade
kommer der dog Sandhed ind i det Bestaaende, det forsvarer
sig ved at dømme sig selv, det vedgaaer den
kristelige Fordring« o.s.v. Men at denne Oplysning ikke
gjør Sageu klar, vil Enhver kunne sige sig selv. Det   % PRINTED AS IS (p. 22): „Sageu“: a turned n
er jo en formelig Afladshandel, den Enkelte her, om
end, hvad Fordelen angaaer, i fuldkommen Uegennyttighed
--- underhaanden og i al Stilhed --- har villet drive
med det Bestaaende.

Om »en verdslig Letsindighed, der i Aaret 1851
lever hen i den Indbildning at være Kristen uden maaske
at have noget Indtryk af Kristendommen«, siger han,
at den kun letsindig og negativt har taget Kristendommen
forfængeligt og er --- »hvis man altid endog tør
dømme den saa strengt« --- til at smile af, »fordi man
% --- p. 23 ---
\setcounter{footnote}{0}%
\opage{23}haaber«. Heri skjuler sig da umiskjendelig den Bagtanke,
at hvis det Bestaaende blot vilde vedgaae sin
»verdslige Letsindighed« og indrømme »Fordringen«, saa
skulde den --- endog uden at dens Tilstaaelse om sin
Afstand kunde (se anførte Artikel i »Fædrelandet«) »kaldes
en Stræben i Retning af at komme Fordringen nærmere«
--- fremdeles have Lov til at kalde sig et kristeligt
Bestaaende og dog blive fritaget for »den egenlige
Kristi Efterfølgelse«. En besynderlig Trafik mellem Ide
og Virkelighed!

I Aaret 1855, da Forfatteren af »Indøvelse i Kristendom«
har forvisset sig om, at det Bestaaende, kristeligt,
er uholdbart, er Naadens Tid omme. Det Bestaaende
lyses i Band, hele Kristenheden belægges med
Interdikt. Endnu i Aaret 1854 kunde det Bestaaende
have frelst sig ved »at trække i anden Potens paa Naaden;
men 1855, da Øieblikket kom, Øieblikket Nr. 1-2-3,
var det forbi med den mulige Aflad, forbi, forbi med
at trække paa Naaden. Saa kunde Virkeligheden folde
sine Hænder, istemme credo, krybe paa Knæerne, klæde
sig i Sæk og Aske, kort, tee sig som den vilde, det
var forbi. Paa Innocens den Tredies Tid, da Frankrig
skulde belægges med Interdikt, optraadte Legaten,
klædt i himmelblaa Stola, som paa Verdens Frelsers
Lidelses Dag, lod Helgenbillederne indhylle i Flor og
Stene kaste fra Prædikestolen ned blandt Menigheden,
til et Vidnesbyrd om, at saaledes havde Himlens og
Jordens Herre nu forkastet det hele Folk. Det var et
gribende Syn, et gruopvækkende Optrin. Men for
alvorligt Tænkende var det dog endnu mere gruopvækkende
at see hin Enkelte med det Nye Testamente
% --- p. 24 ---
\setcounter{footnote}{0}%
\opage{24}i Haanden, en Poeniterende, der gik ind til Dommen,
en Døende, allerede klædt i Ligskjorten, kaste Stene paa
den gamle Kirke og bryde Staven over den hele bestaaende
»Kristenheds Kristendom.« Det skete forunderlig
brat og dog med en Evigheds Alvor; det var som
en Overilelse og dog saa vel betænkt, saa længe, saa
alsidig forberedt. Et frygteligt Brud mellem Ide og Virkelighed!

Slige Spring fra den ene Yderlighed til den anden,
lade sig kun forklare af en voldsom Spænding, en indre
Strid i Principet selv. Først seer det ud, som var Ideen,
hvad alt det Udvortes angaaer, fuldkommen tilfreds med
den givne Virkelighed. Ideen søger den Enkelte. Den
Enkelte har nok i sig selv, thi han har nok i Ideen.
Det Eneste, han for Inderlighedens Skyld trænger til, er
nogle høist gavnlige Øvelser i Taalmod, Selvfornægtelse,
Overbærenhed, Langmodighed, Fredsommelighed o. s. v.,
men i denne Hensende vil den bestaaende Virkelighed   % PRINTED AS IS (p. 24): „Hensende“ for Henseende
til enhver Tid kunne være ham til Tjeneste; thi det er
just det Gode ved Virkeligheden: den lader det sjeldent
mangle paa en eller anden skjøn Leilighed til at indøve
hine for den Enkeltes Frihed og Fred uskatterlige Dyder.
Paa denne Vis er Ideen mod den ydre Omgivelse
da Mildheden, Føieligheden, Nedladenheden selv. Ingen
Ophævelse for den Enkeltes Skyld, ingen Uleilighed
med de udvortes Betingelser; Alt er godt, meget godt,
udmærket godt, som det er. Men hvad betyder dette?
Det betyder, at Ideen i absolut Inderlighed er uendelig
sikker paa sig selv, men det betyder ingenlunde, at
den er i god Forstaaelse med Virkeligheden. Langtfra!
det betyder, at den Enkelte med dulgt Ringeagt for alt
% --- p. 25 ---
\setcounter{footnote}{0}%
\opage{25}Udvortes seer smilende ned paa det Bestaaende, i den
stiltiende Forudsætning, at det Bestaaende med taus Beundring
ser op til ham, og ydmyg erkjender Ideens Overlegenhed.
Da nu Virkeligheden desværre ikke altid er
saa fintfølende, og det Bestaaende helst gaaer efter sit
eget Hoved, saa kan det ved mindste uforudselige Anledning
skee, at en voldsom Krig bryder ud. Thi Ideen
gjør store Fordringer, det Ubetingede vil herske. Saa
huldsalig mild som den naadige Ide ellers kan være,
saa forfærdelig streng er igjen den vrede Guddom. Var
Alt forhen Ideen tilpas, saa er nu Intet tilpas: Virkeligheden
er dømt, der holdes Standret over det Bestaaende.
Kampen føres paa ulige Vilkaar. Det Bestaaende
har til sit Værge fysisk Magt; den Enkelte,
Ideens Tjener, strider kun med aandelige Vaaben. Det
er Ideens Forherligelse, dens høieste Aabenbarelse, at
dens sande Stridsmænd kunne seire indvortes ved at
tabe udvortes, at flere og flere Enkelte kunne blive Martyrer.
Saalænge kan Striden fortsættes, saamange Offere
kunne falde, at den plumpe, brutale Virkelighed
tilsidst for sin egen Skyld bliver nødt til at krybe til
Korset og gjøre Ideen Indrømmelser, paa hvilken Betingelse
der da kan sluttes, om ikke en egenlig Fred,
saa dog en Vaabenstilstand. Men enten nu Tilstanden
forøvrigt er fredelig eller ufredelig, saa er Hemmeligheden
kommen for Lyset. Forholdet imellem Ide og
Virkelighed er nemlig bestemt saaledes, at det Ubetingede
maa ligge i en evig Strid mod sine Betingelser.
I Hemmelighedens Indre sidder en Modsigelse.
Modsigelsen er tveegget: den gjennemskærer Virkeligheden,
og den saarer Ideen.

% --- p. 26 ---
\setcounter{footnote}{0}%
\opage{26}Fra Virkelighedens Side er Modsigelsen den, at Virkeligheden
med sit Væv af Betingelser idelig hungrer
efter Ideen og dog ikke kan fordøie Ideen, efterdi de
endelige Betingelser aldrig kunne forliges med det Ubetingede.
I Ideen selv er Modsigelsen den, at den vil
virkeliggjøres, at den inderlig begjærer Virkelighed uden
dog at kunne fordrage Virkeligheden, efterdi det Ubetingede
bestandig frastødes af Betingelserne. Skal en
sand gjensidig Forsoning mellem Ide og Virkelighed være
mulig, da maa Principet selv forandres.

Det er en Grundenhed af det Ubetingede med dets
Betingelser, hvorpaa Alt væsenlig kommer an. Der maa
i Ideen være Bestemmelser af det Betingede, der passe for
den ydre Virkelighed, og der maae i Virkelighedens eget
Væsen gives Bestemmelser af det Ubetingede, som idetmindste
tilnærmelsesvis --- altsaa under en fortsat Udvikling
--- tilfredstille Ideen. Spørgsmaalet er nu, hvorledes
det forholder sig med disse vexelsidige Bestemmelser.

Hvad S. K. angaaer, saa har det hidtil været og
er endnu vor fulde Overbevisning, at han med sit skarpe
Blik meget vel saae de forsonende Mellembestemmelser,
men at han ifølge den til hans eiendommelige Opgave
svarende Taktik, ansaae sig forpligtet til at holde dem
i Baggrunden. Han maatte holde strengt paa det Ubetingede
og see at virke ved et pludseligt Slag, dersom
det ellers skulde lykkes ham at vække Samtiden op af
dens aandelige Døs og anbringe et Afgjørende. »Thi«,
siger han, (Øieblikket Nr. 1) »at anbringe et Afgjørende
--- og dette er Opgaven --- lader sig ikke gjøre
paa samme Maade som alt Andet; og naar nu tilmed
Tidens Ulykke er dette »til en vis Grad«, til en vis
% --- p. 27 ---
\setcounter{footnote}{0}%
\opage{27}Grad at gaae ind paa Alt, naar dette just er Sygdommen,
saa maa der for Alt passes paa, at det saavidt
det er muligt ikke skeer, at den ogsaa til en vis Grad
gaaer ind herpaa, hvorved Alt er tabt. Nei, et Afgjørende
anbringes anderledes end Andet. Som Rovdyrets
Spring paa sit Bytte, som Ørnens Stød i Faldet, saaledes
anbringes et Afgjørende: pludseligt og samlet i
Et (intensivt). Og som Rovdyret ener Snildhed med
Kraft, først snildt holder sig ganske stille --- saaledes
stille kan intet tamt Dyr være; og saa samler sig helt
i et eneste Spring eller Stød --- saaledes kan intet
tamt Dyr samle sig til eller reise sig i Springet: saaledes
anbringes et Afgjørende. Først stille --- saa stille er
Veiret aldrig paa en stille Dag, saa stille er det kun før
Torden; og saa bryder det løs«.

Den Grundform, der alene er istand til at optage
Mellembestemmelserne af det Betingede og det Ubetingede
og forlige Ideen med Virkeligheden, er med sit
rige Indhold mere end en Form, den er selv en Ide,
en sand Virkelighedside, nemlig Samfundets Ide. Det
er Samfundsideen, som S. K. overalt stiller i Skygge,
for at der kan falde et desto stærkere Lys paa hans Enkelte.
»Den Enkelte«, siger han, »er den Kategori, gjennem
hvilken, i religiøs Henseende, Tiden, Slægten, Historien
skal. Og han, der stod ved Thermopylæ, var
ikke saaledes sikkret som jeg, der for dog idetmindste
at faae gjort opmærksom paa den, har staaet ved denne
Snævring »den Enkelte«. Han skulde nemlig forhindre
Skarerne fra at trænge igjennem Snævringen; trængte
de igjennem, saa havde han tabt. Min Opgave er idetmindste
langt mindre udsættende mig for den Fare at
% --- p. 28 ---
\setcounter{footnote}{0}%
\opage{28}blive nedtraadt, da den var som ringe Tjenende\dots{} om
mulig at foranledige, at indbyde, at bevæge de Mange
til at trænge igjennem denne Snævring »den Enkelte«,
gjennem hvilken dog, vel at mærke, ingen trænger uden
ved at blive den Enkelte; det Modsatte er jo en kategorisk
Umulighed.« Her antages altsaa, at Snævringen
»den Enkelte« er ganske anderledes sikker end Snævringen
ved Thermopylæ; men derom maa det være tilladt
at nære en Tvivl.

Snævringen ved Thermopylæ havde dog altid, hvad
enten Skarerne trængte igjennem eller ikke, sin bestemte
Strækning, sine klippefaste Grændser; men Snævringen
»den Enkelte« har, saavidt vi kunne see, ingen
klippefaste Grændser: den udvider sig bestandig, idet
Skarerne trænge igjennem, endog i den Grad, at den
tilsidst maa rumme dem alle. Thi, naar Alle blive de Enkelte,
saa maae de Enkelte nødvendigvis igjen blive Alle,
men at blive Alle er jo at blive Mængden: ved at gaae
ud af Mængden --- ind i Snævringen, gaae Skarerne
altsaa netop ind i Mængden --- ud af Snævringen.
For at indgaae i Sandheden maae de Enkelte uden Undtagelse
udskille sig fra Mængden, thi Mængden er Usandheden;
men naar nu hver Enkelt, efter at de Alle have
udskilt sig fra Mængden, just af den Grund kommer til
at staae midt i Mængden: hvor befinder han sig saa, i
Sandheden eller i Usandheden? Sige vi i Usandheden,
efterdi han er En af Mængden, og Mængden er Usandheden,
saa er her i Snævringen begaaet en Feil, her er,
rentud sagt, skeet en Ulykke. Naar der gaaer Ild i en
stor Bygning, hvor Mange ere forsamlede, og Bygningens
eneste Udgang er en Snævring, hændes det let, at
% --- p. 29 ---
\setcounter{footnote}{0}%
\opage{29}Alle brænde inde, fordi Alle paa engang ville ud: ved
at gaae enkeltvis, En efter En, kunde dog idetmindste
Nogle været frelste. Omtrent saaledes forholder det sig
med Snævringen »den Enkelte«, naar det skal gjælde
ubetinget, at Mængde er Usandhed. Kun ved at være
Undtagelsen, den fremfor de Andre Udmærkede, kan den
Enkelte kategorisk adskille sig fra Mængden. Hver
Slægt, hver Skare maa man altsaa, naar den Enkelte
skal frem, tilraabe: En ad Gangen, En i hver Slægt,
kun En, en Eneste! thi En til og En til er allerede
Mængde.

Det er øiensynligt, at Mængden umulig kan danne,
hvad S. K. kalder en \emph{kvalitativ} Modsætning til den
Enkelte. Mængden dannes af Enkelte; som de Enkelte
ere, saaledes er Mængden. Bestaaer Mængden af sande
Enkelte, saa er der Sandhed i Mængden; er Mængden
Usandhed, saa er Grunden den, at den bestaaer af
usande Enkelte. Mængden som Mængde er i og for sig
hverken god eller ond, hverken sand eller usand. For
at opfatte Mængden som den tilværende Usandhed, maa
man først i sit Begreb udelukke den fra Sandheden, ved
at udelukke den fra Ideen. Den ideløse Mængde --- ja
det er uimodsigelig sandt --- den ideløse Mængde er efter
Omstændighederne Masse, Pøbel, Publikum, det Numeriske,
det Usande; men den er da ikke det Usande,
fordi den er Mængde, den er det kun, fordi den er udelukket
fra Sandheden. Paa tilsvarende Maade kan man
jo ogsaa begynde med at udelukke den Enkelte fra
Sandheden, betragte de Enkelte udenfor Ideen; og det
følger da ganske af sig selv, at de ideløse, det vil sige
usande, Enkelte ere usande. Men de Enkelte ere da
% --- p. 30 ---
\setcounter{footnote}{0}%
\opage{30}usande, ikke fordi de ere Enkelte, men fordi de ved
det opstillede Begreb ere udelukkede fra Sandheden.

Skal der dannes nogen sand Modsætning mellem
Mængden og den Enkelte, da maae begge de modsatte
Led falde indenfor Ideen; men Mængden, naar den henføres
til sin Ide, er det tilværende Samfund. Saaledes
komme vi da til Modsætningen mellem de Enkelte og
Samfundet; Mængde er Virkelighed, Samfund Ide: vi
maae da særlig lægge Vægt paa det Spørgsmaal: hvorledes
det i det virkelige Samfund forholder sig med Ide
og Virkelighed.

Samfundet har sin Bestaaen i den organiserede
Mængde. Naar nu Ideen skal udvikle og virkeliggjøre
sig i Samfundet, hvilken Vei gaaer den saa? Begynder
den delvis med de Enkelte og arbeider sig først lidt
efter lidt igjennem det Hele, eller begynder den strax
ud fra Helheden, for efterhaanden at tænde sit Lys i
de Enkelte, eller begynder den paa een Gang baade ud
fra Helheden og ud fra de Enkelte? Paa disse for Frihedsideen
høist vigtige Spørgsmaal giver S. Kierkegaard
os ingensteds noget bestemt Svar. Han skyder Undersøgelsen
fra sig, forudsætter uden videre Samfundets
Bestaaen, polemiserer mod Mængden og holder paa den
Enkelte. Men at Undersøgelsen netop for den Enkeltes
Skyld er af yderste Vigtighed, vil blive indlysende, naar
vi kaste et Blik paa de Enkeltes Stilling i de forskjelligartede,
snævrere og videre Samfundskredse, i Menigheden,
i Staten, i Folket.

Stille vi, for at begynde med det Religiøse, den Enkelte
mod Menigheden, saa møde vi strax den Kjendsgjerning,
at saalænge Sandheden udelukkende var til
% --- p. 31 ---
\setcounter{footnote}{0}%
\opage{31}som Forbilledet selv, den Ene, den ubetinget Enkelte,
kunde Menigheden ikke stiftes, og da Menigheden var
stiftet, kunde ingen Enkelt træde selvstændig op imod Menighedens
Aand. Men af Iver for at forsvare Snævringen
»den Enkelte« gaaer S. K. i denne Anledning saa vidt,
at han endog paastaaer, at »Kristendommen egenligen
ikke er kommen ind i Verden«, da end ikke Apostlene
ret have formaaet at holde Orden ved Indgangen. »Thi
Et er,« siger han (Øieblikket Nr. 5), »at virke saaledes
for Udbredelse, at man uafbrudt tidlig og sildig forkynder
Læren for Alle, noget Andet er, at være for hurtig
til at tillade Folk saadan i 100, i 1000-vis at antage
Navn af Kristne\dots{}. Om en Folke-Forsamling havde
ladet sig gribe af Kristi Tale, derfor havde han vistnok
ikke strax tilladt disse Tusinder at kalde sig Kristi
Disciple. Nei, han holdt stærkere igjen. Han vandt
derfor ogsaa i 3½ Aar kun 11 --- medens en Apostel
paa \emph{een} Dag vel i \emph{een} Time vinder 3000 Kristi Disciple.
Enten er her Disciplen større end Mesteren, eller Sandheden
er, at Apostelen er lidt for hurtig til at slaae
til« o. s. v.

Skal dette tages efter Bogstaven, saa er Menighedens
Ide uden videre at opgive. Menighedsideen forudsætter
nemlig, at Aanden paa en Gang begynder sin
Virksomhed ud fra Samfundet og ud fra de Enkelte.
Apostelen taler som den Enkelte til de Enkelte i Mængden;
men han taler af Aanden. Den Aand, hvoraf han
taler, er den religiøse Samfundsaand, den Aand, der
ikke vil virke enkeltvis paa de spredte Enkelte, men
forene dem i Samfundet, for i og ved Samfundet at
virke i enhver Enkelt derved, at de kunne virke paa
% --- p. 32 ---
\setcounter{footnote}{0}%
\opage{32}hverandre indbyrdes og under indbyrdes Vexelvirkning
udvikle sig til sande Enkelte. Derfor iler han med at
faae Menigheden stiftet. Skulde Apostelen have ventet.   % PRINTED AS IS (p. 32): „ventet. indtil“: a stop before a lower-case word
indtil de Enkelte vare blevne mere forberedte, for at
Sandheden i Menigheden omsider kunde blive et slet og
ret Resultat af Sandheden i de udenfor Menigheden
dannede Enkelte, da var han aldrig kommen til at stifte
nogen Menighed. Naar skulde han vel have stiftet den?
Før de Enkelte endnu vare blevne sande Enkelte? Det
var for tidlig! Efter at de vare blevne sande Enkelte?
Det var for sildig! En anden Sag er det, hvis Apostelen
havde ladet det blive ved at døbe de mange Tusinde
og sagt til Mængden: nu ere I døbte, nu ere I Menigheden,
nu behøves der Intet videre; men saaledes gik
det jo ikke til. Det er saa langt fra, at Mængden som
Mængde kommer til at raade i Menigheden, at Menighedens
Aand nu netop sætter sine Enkelte, sine Udvalgte,
den store Opgave at tale Mængden til Rette.   % PRINTED AS IS (p. 32): „Rette. organisere“: a stop before a lower-case word
organisere, tugte, opdrage den, at de Mange, hver især,
maa erkjende, hvad det vil sige at være for Gud den
Enkelte.

Stille vi, hvad det Politiske angaaer, den Enkelte
imod Staten, da udkommer et lignende Resulat. Massen,   % PRINTED AS IS (p. 32): „Resulat“
den ideløse Mængde kan naturligvis ikke stifte nogen
Stat; men den Enkelte, der foragter Mængden, kan
det heller ikke. Statens Ide har to lige oprindelige
Udgangspunkter. Den rører sig allerede kraftig i Mængden,
idet den begynder at klare sig for de enkelte
Ledere. For at organisere Masserne og modarbeide
Publikum, det blot Numeriske, er det ikke nok »at foranledige,
at indbyde, at bevæge de Mange til at trænge
% --- p. 33 ---
\setcounter{footnote}{0}%
\opage{33}igjennem Snævringen den Enkelte«, det er tillige nødvendigt,
at man uafbrudt tidlig og sildig indskærper
Samfundsideens Gyldighed, Regeringens Myndighed, Borgerens
Rettigheder og Borgerens Pligter. Ved at skille
Bevidstheden om de borgerlige Pligter fra den religiøse
Inderlighed og indprente de Enkelte Ringeagt for Samfundets
Opgaver under Foregivende af, at det kun er
verdslige udvortes Former, gjør man Staten en maadelig
Tjeneste. Hvad kan det hjælpe at bukke for Autoriteterne
og tale fromt, underdanigt om, at Øvrigheden er
indsat af Gud, og at Styrelsen bør overlades »de ordenligen
dertil Beskikkede og Kaldede«, naar man i samme
Øieblik underminerer Grundlaget for al fornuftig Regering?
Hvem er det, der skal regere? de Enkelte eller
ikke de Enkelte? Dersom det ikke er de Enkelte, saa
maa det jo være Mængden, Publikum, det Numeriske.
Men dersom det kun er »Mængdens Ledere«, der udgjøre
Regeringen, da klinger det jo besynderligt, at en
saadan Regering ikke væsenlig bør søge sin Styrke i
»verdslig Klogskab«, men i »en inderlig Forvandling i
Retning af at blive fast ved at frygte Gud«. Skal Regeringen
blive fast ved at frygte Gud, da maa den nødvendigvis
bestaae af ordenligen dertil beskikkede Enkelte,
af Saadanne, som i religiøs Inderlighed have lært,
hvad det er at trænge igjennem »Snævringen«. I disse
sande Enkelte kan det Religiøse umulig forholde sig
ligegyldigt til det Politiske; for at faae Styrke til at
regere, maae de blive faste ved at frygte Gud, men den
Inderlighedskraft, de vinde ved at frygte Gud, skal igjen
lægge sig for Dagen i den Kraft og Myndighed, hvormed
de regere: en inderligere Vexelvirkning mellem Politik
% --- p. 34 ---
\setcounter{footnote}{0}%
\opage{34}og »Gudsfrygt« lader sig ikke tænke. Hvad der gjælder
om nogle sande Enkelte, gjælder om alle; de skulle
jo alle igjennem Snævringen. I Egenskab af Statsborgere
have Samfundets Enkelte nu det skjønneste
Exempel i Regeringen. Ligesom den »gudfrygtige« Regering
fra sin Side paa ingen Maade tør ansee de ydre
sociale Former for ligegyldige, uagtet de ere verdslige,
saaledes tør de gudfrygtige Borgere fra deres Side det
ikke heller. Er det en Samvittighedssag for en af sande
Enkelte sammensat Regering at fremme Samfundsideens
Udvikling og sørge for Statens Vel ved at vaage over
Borgernes Rettigheder, saa maa det ogsaa være en Samvittighedssag for alle enkelte Borgere at leve sig ind i
Samfundets Ide og punktlig opfylde de borgerlige Pligter.

Endskjøndt den religiøse og politiske Ide vel maae
modsættes hinanden, som den indre Frihed modsættes
den ydre, saa have de dog det Fælles, at de begge ere
Samfundsideer, begge anlagte paa at indlade sig med
Mængden, optage Betingelserne i det Ubetingede og omdanne
Virkeligheden. Den Rolle, Mængden spiller ved
Bestemmelsen af Forholdet mellem Ide og Virkelighed,
viser sig især fremtrædende, naar Talen er om Samfundet
i videste Omfang, om \emph{Folket.}

At \emph{Folket} i et Statssamfund maa, naar vi see paa
Virkeligheden, være til som Mængde, kan ikke nægtes;
men just af den Grund bliver det for Ideens Skyld
af største Vigtighed at gjøre Forskjel paa Folket og
Mængden. Folket er Mængden, forsaavidt det er Summen
af Befolkningens Individer, og dog er Mængden endnu
ikke Folket; kun ved at henføres til sin oprindelige Ide,
Ideen af det Grundmenneskelige i naturkraftig Eiendom%
% --- p. 35 ---
\setcounter{footnote}{0}%
\opage{35}melighed (Jvnfr. »Om Hindringer og Betingelser for det
aandelige Liv i Nutiden«, VIII), kan Mængden opfattes
som Folk. Var det virkelig umuligt at gjøre Forskjellen
mellem Folket og Mængden gjældende paa en indtrængende,
klar, for den sunde Fornuft indlysende Maade,
saa vilde Folkefrihed og Mængdefrihed være enstydige
Ord. Men dersom Folkefrihed uden videre falder sammen
med Mængdefrihed, saa at en fri Forfatning skal
være den, der opstiller Betingelser for, at Mængden paa
lovlig Maade kan faae sin Villie, om denne Villie end
i sig selv er nok saa ufornuftig, sat igjennem: da er det
ikke til at indsee, hvorledes noget fornuftigt Menneske
skal kunne sætte enten Lid eller Haab til en saadan
Frihed. Det er kun Folkefrihedens Modstandere, som,
hvad enten de forøvrigt foragte eller forgude Mængden,
gjøre Folket til Et med Mængden.

I hvilken Betydning Folket som Folk maa siges at
være hævet over sin egen Mængde og den sande Folkevillie
ophøiet over den tilfældige Flertalsvillie, kan let oplyses ved
nogle Exempler. Sæt, at Flertallet i en Befolkning, hvad der
vel er lidet sandsynligt, men dog ikke utænkeligt, i den Grad
hengav sig til Troen paa den evige Fred, at det ansaae
alt Forsvarsvæsen for en Overflødighed, og i den Anledning
fik sat igjennem, at alle Udgiftsposter til Land- og
Søværn bleve strøgne: saa var det dog ikke Folket,
der nærede en saadan Tro, det var kun Mængden.
Eller sæt, at en Afstemning engang med stor Majoritet
udvirkede Tilladelse til, at Mandfolk og Fruentimmer ---
for de personlige Friheders Skyld --- herefter maatte
klæde sig ens, saa at Kvinden, for at komme til sin
naturlige Ret, endog turde anlægge forlorent Skæg, saa
% --- p. 36 ---
\setcounter{footnote}{0}%
\opage{36}var det dog ingenlunde Folket, der kom tilorde ved en
saadan Afstemning, nei, det var kun Mængden. Eller
sæt, at et frisindet Flertal, i Erkjendelse af, at Ægteskabet
er en ligesaa borgerlig som naturlig Indretning,
fik ved en syv og tyvende Behandling af Ægteskabsloven   % PRINTER'S DEFECT (p. 36): „fik“: the f did not ink its head (both OCRs read t); letter supplied
den Forandring sat igjennem, at det for Fremtiden
skulde være en Mand tilladt at have flere Koner,
kun at Konernes Antal kom til at staae i strengt normeret
Forhold til Indkomstskatten, saa vilde et Fremskridt af
denne Natur dog ikke være at tilskrive Folket, men
Mængden. Hvad er nu det, som i de anførte Exempler
gjør det saa let, at adskille Mængdens Griller fra Folkets
Tanker? Det er jo netop det i sig selv Urimelige,
Ufornuftige, for det Hele Skadelige; thi Folket som Folk
kan kun ville det Fornuftige, det med Livets Grundfordringer
Stemmende, det for hele Samfundet Gavnlige.

Dersom alle for Samfundet vigtige Anliggender vare
af den Natur, at der ikke behøvedes Andet end sund
Sands, naturlig Forstand og en god Villie for strax at
indsee, hvad det er, hvorpaa det hver Gang kommer an,
og hvormeget der i kritiske Øieblikke kan staae paa Spil,
saa vilde det endda være en let Sag at klare Folkets
Syn og tale Mængden tilrette. Men Vanskeligheden er,
at netop de allerbetydeligste Opgaver, de store, i Samfundslivet
dybest indgribende Spørgsmaal, de i egenligste
Forstand folkelige Anliggender høre til dem,
hvis Betydning vanskeligst fattes og tilfulde erkjendes
af det hele Folk, fordi Mængden, som det hedder, ikke
kan følge med. Skal denne for sand Folkedannelse hidtil
uoverstigelige Hindring omsider overvindes, da kan
det kun skee paa den Betingelse, at de folkelige
% --- p. 37 ---
\setcounter{footnote}{0}%
\opage{37}Spørgsmaal paa en Gang behandles baade theoretisk
som Idespørgsmaal og praktisk som Virkelighedsspørgsmaal:
Ide og Virkelighed maae da ikke skilles ad, men
holdes mod hinanden, at det i klare enfoldige Træk kan
sees, hvorledes de ved enhver paagjældende Sag belyse
hinanden gjensidig.

Spørgsmaalet om Folkereligionen er, hvad vist Enhver
vil indrømme, baade et Idespørgsmaal og et Virkelighedsspørgsmaal.
Den almindelige, i sig selv meget
rigtige Mening, at et Folk umulig kan undvære en positiv
Religion, er for svævende og svag til at bære Virkelighedens
tunge, dybtskærende Forhandlinger angaaende
Folkekirkens Stilling til Staten. Det første Fornødne
er, at Folket i sin Helhed kommer til Klarhed over Religionens
Væsen og indseer dens Betydning for Menneskelivet.
Først efterhaanden som en levende Erkjendelse
af, hvad Religion er, i hvilket indre væsenligt Forhold
den staaer til ethvert selvstændigt Folks Sæder og Tænkemaade,
til Kulturens og Civilisationens Fremgang i alle
Retninger, mere og mere udbreder sig, vil Folket optræde
som Folk og lægge sin Stemnings og Stemmes
Vægt paa Forhandlingernes Vægtskaal.

Hvilke Ordfordreielser og Misforstaaelser, der kunne
sættes i Omløb, naar store Idespørgsmaal uden videre
trækkes ned i Endeligheden og afgive Slagord til Bedste
for Partierne, sees af Nutidens besynderlige Strid om
Intelligensens Værd. Udenfor Ideen vil en saadan Strid
aldrig ende. Det er aabenbart, at de, der kaste Skygge
paa Intelligensen, forbinde med Ordet en ganske anden
Forestilling end de, der forsvare den. Intelligens er
Indsigt, Indsigt kræver Oplysning. Uden Oplysning in%
% --- p. 38 ---
\setcounter{footnote}{0}%
\opage{38}gen Frihed, uden folkelig Oplysning ingen Folkefrihed.
Hvad er vel Frihed uden Oplysning, et frit Folk uden
Intelligens? En Dag uden Lys, et Sværd uden Klinge, ---
tankeløse Modsigelser! Mod den indbildte Dannelse, den
tomme Oplysning, den opblæste, ufrugtbare Intelligens
bør der gjøres Indsigelse; men begrundede Indsigelser
mod den falske Intelligens indeholde jo netop Forsvar
for den sande.

Paa Virkelighedens Overflade kan det vel i Mængdens
Øine undertiden se ud, som om Konst, Poesi og
Videnskab kun vare Luxusartikler til Nydelse for de
høiere Samfundsklasser; men et Folk, der vil forstaae
sig selv, vil ogsaa forstaae, at Aandens klare Former
høre væsenlig med til Folkelivets frie, alsidige Udvikling,
ere sande Fornødenheder, endskjøndt de ikke, som
»\emph{Brød og Jern}« ligefrem tør regnes til det første
Nødvendige. Grundspørgsmaalet er, hvorledes den
høiere Dannelse i det Hele forholder sig til den lavere;
thi at den lavere ikke kan hæve sig uden Hjælp fra den
høiere, er en Kjendsgjerning. Spørg Arbeideren, han
veed allerede Besked, at der kun tilgavns kan arbeides
med Haanden, naar der tillige arbeides i Aanden, at
han først tilfulde kan bruge Kræfterne, naar han har
lært at bruge Tankerne.

Et talende Bevis for, at Oplysning og Frihed gaae Haand
i Haand, er den umiskjendelige Trang, den Drift og Stræben,
der overalt rører sig i Befolkningen efter Kundskab
og Dannelse: Masserne samle sig, hvor de ane en Tanke;
de lytte til Ordet, de skimte Lyset, de trænge sig frem,
at hilse Frihedens Dag. Denne Stræben, denne Trang
er et Vidnesbyrd om, at Folket som Folk har en dyb
% --- p. 39 ---
\setcounter{footnote}{0}%
\opage{39}Fornemmelse af sin Fremtids Opgave, at det ikke vil
slaae sig til Ro ved den udvortes Frihed, men skride
videre frem og forstaae at tilkæmpe sig den indvortes.
Først ved Forening af den udvortes Frihed med den indvortes
bliver den folkelige Frihed fuldstændig, betinget i
Virkeligheden, ubetinget i Ideen: Folkets Frihed er, som
Folket selv, en vordende, voxende Enhed af Ide og
Virkelighed.
\signature{R. \emph{Nielsen}.}

\end{document}
